\chapter{本体}
\dirtree{%
.1 本体：对某个领域的概念体系做形式化建模.
.2 目标：让人和机器都能明确理解.
.3 有哪些东西.
.3 怎么分类.
.3 怎么关联.
.3 能推出什么.
.2 0 基础思想层：理解本体的地基.
.3 概念化 Conceptualization.
.4 先决定如何看待这个领域：哪些东西重要，哪些关系重要.
.3 显式化 Explicit.
.4 把隐含理解写清楚，避免“大家心里都懂但机器不懂”.
.3 形式化 Formal.
.4 用机器可处理的语言表达概念、关系和约束.
.3 共享性 Shared.
.4 本体不是个人笔记，而是领域共同体可复用、可对齐的知识结构.
.2 能力一：定义这个领域里“有什么” $\rightarrow$ 概念建模.
.3 类 Class / Concept.
.4 例如：人、学生、课程、疾病、药物.
.3 实例 Individual / Instance.
.4 例如：张三 是 学生；阿司匹林 是 药物.
.3 层级 Taxonomy / is-a.
.4 例如：本科生 是 学生；学生 是 人.
.2 能力二：定义这些东西“怎么关联” $\rightarrow$ 关系建模.
.3 对象属性 Object Property.
.4 连接两个对象：学生 选修 课程；药物 治疗 疾病.
.3 数据属性 Data Property.
.4 连接对象和数值：学生 有年龄；课程 有学分.
.3 关系约束 Constraint.
.4 例如：每门课至少有一位老师；一个人只能有一个出生日期.
.2 能力三：表达规则并自动推出新知识 $\rightarrow$ 逻辑推理.
.3 继承推理.
.4 已知：本科生 是 学生，张三 是 本科生.
.4 推出：张三 是 学生.
.3 属性推理.
.4 已知：hasParent 的反关系是 hasChild.
.4 若 A hasParent B，则推出 B hasChild A.
.3 一致性检查.
.4 若“学生”和“老师”互斥，而张三同时是二者，则发现冲突.
.3 分类推理.
.4 根据条件自动判断某个实例属于哪个类.
.2 能力四：让知识可交换、可查询、可复用 $\rightarrow$ 语义互操作.
.3 RDF.
.4 用三元组表达事实：主体 - 谓词 - 客体.
.3 RDFS.
.4 给 RDF 加上类、子类、属性层级等基本语义.
.3 OWL.
.4 更强的本体语言，可表达复杂约束并支持推理.
.3 SPARQL.
.4 查询语义数据，就像 SQL 查询关系数据库.
.2 能力五：用本体解决实际问题 $\rightarrow$ 应用层.
.3 知识图谱建模.
.4 本体规定 schema，知识图谱填入具体事实.
.3 数据融合.
.4 把不同来源的数据按统一语义对齐.
.3 智能问答.
.4 让系统理解“人、机构、地点、事件”之间的语义关系.
.3 领域知识库.
.4 医疗、法律、金融、工业、教育等领域的知识标准化.
.3 AI/RAG 增强.
.4 用本体约束概念、关系和推理路径，让检索和生成更可靠.
}
\begin{definition}[Ontology]
An ontology is a formal, explicit specification of a shared conceptualization.本体是对共享概念化的形式化、显式规约
\end{definition}
\begin{dxtips}
    formal 表示机器可处理，explicit 表示概念和约束被清楚定义，shared 表示不是个人理解，而是某个共同体认可的领域知识。
\end{dxtips}
    
\begin{dxtips}
    本体其实就是像在特定领域构建出数学这种公理定义定理体现或者像英语语法的主谓宾定状补规则结构
\end{dxtips}
\[
\begin{array}{c c}
\textbf{数学体系} & \textbf{本体体系} \\
\hline
\text{定义} & \text{类 / 概念} \\
\text{对象} & \text{实例} \\
\text{关系} & \text{属性 / 关系} \\
\text{公理} & \text{约束 / 逻辑规则} \\
\text{定理} & \text{推理出来的新知识} \\
\text{证明} & \text{推理过程}
\end{array}
\]